\documentclass[11pt]{letter}
\usepackage[left=2.5cm,right=2.5cm,top=2.5cm,bottom=2.5cm]{geometry}
\usepackage{hyperref}
\usepackage{charter}
\usepackage[T1]{fontenc}
\usepackage[utf8]{inputenc}

% letter 클래스에는 \paragraph 가 없다 — 소제목용으로 정의
\makeatletter
\providecommand{\paragraph}{}
\renewcommand{\paragraph}[1]{\par\medskip\noindent\textbf{#1}\par\nobreak\smallskip}
\makeatother

\signature{Ku Kang, Ph.D. \\
CBRN Defense Research Institute \\
Ministry of National Defense, Republic of Korea \\
E-mail: bisu9082@gmail.com}

\address{Ku Kang \\
CBRN Defense Research Institute \\
Seoul 06796, Republic of Korea}

\begin{document}

\begin{letter}{The Editor \\
\textit{Chemical Science} \\
The Royal Society of Chemistry \\
Burlington House, London W1J 0BA, UK}

\opening{Dear Editor,}

We submit for your consideration a manuscript entitled
\textbf{``A metal chelate's effect on molecular recognition is not a property of
the chelate''} for publication in \textit{Chemical Science} as an Edge Article.

\paragraph{The assumption we tested, and why it matters beyond one field.}
Coordination chemistry is routinely appended to molecular recognition: a metal
complex is attached to a molecule that binds a biological target, and the two
jobs---holding the metal and still binding---are treated as separable. That
separability is assumed rather than tested, and it is load-bearing: chelators
are selected from published comparisons and from heuristic metal--chelator
matching tables. We audited the assumption against the literature's own data and
found it does not hold. The effect of the chelate on recognition is real and
large, but it is not encoded in the chelate.

\paragraph{What the audit shows.}
From ten studies we assembled 106 matched affinity measurements in which a
single targeting molecule was conjugated to two or more chelators and assayed
under one protocol, giving 32 matched sets across five target families. Within a
set, chelator identity alone shifts affinity by a median factor of 2.6 and by up
to 25.3, while propagated measurement error accounts for only 11.1\% of the
within-set variance. None of that signal transfers: under
leave-one-target-family-out validation an empirical chelator ranking
($R^2 = -0.250$), classical physicochemical descriptors ($R^2 = -0.219$) and
solvated GFN2-xTB descriptors computed on 27 structurally verified complexes
($R^2 = -0.340$) all perform worse than assuming no effect, none orders a
held-out set better than chance, and a within-set permutation test places every
model inside the no-information null. What the data do contain is interaction:
chelator identity crossed with target family explains $R^2 = 0.498$ against
$R^2 = 0.133$ for the chelator alone.

\paragraph{Three features we expect your referees to test, and our answers.}
\textit{(i) ``The quantum chemistry was not good enough.''} We anticipated this
and closed it. Complexation energies referenced to the bare ion are dominated by
Coulombic neutralisation---they span $-691$ to $-2257$\,kcal\,mol$^{-1}$ and
track the metal's formal charge at $r = -0.914$---so we replaced them with a
ligand-exchange energy in which the ion term cancels exactly
($r = -0.199$, $p = 0.320$), recomputed every complex with ALPB implicit
solvation and wider conformer sampling, and verified coordination mode against
the known mode for all 27 metal complexes. Improving the physics raised
in-sample explanatory power ($R^2 = 0.070 \rightarrow 0.139$) and \emph{lowered}
held-out performance ($R^2 = -0.227 \rightarrow -0.340$). The negative result
survives its own strongest objection.
\textit{(ii) ``The sample is small.''} We bound rather than excuse this: an
injection analysis shows the design would detect a transferable twofold effect
about three times in four, and is blind below $\sim$1.4-fold.
\textit{(iii) ``It is a lipophilicity artefact.''} Every log$D_{7.4}$ or
log$P$ value reported alongside the affinity data was transcribed and tested; it
explains 7.4\% of the within-set difference and is not significant.

\paragraph{Why \textit{Chemical Science}.}
The claim is a chemical one of general interest: a property of an isolated
complex cannot encode a contact the complex never makes, and we show this
quantitatively rather than asserting it. The work also belongs to a line your
journal has recently championed---the auditing of widely used computational
evaluations, as in Schuh \textit{et al.}, \textit{Chem. Sci.}, 2026, DOI:
10.1039/d6sc01799a, which exposed hidden leakage and redundancy in biomolecular
benchmarks. Our audit shares that methodology---leakage-resistant splits,
controlled injection experiments, permutation nulls---and applies it to a
chemical design rule rather than to a benchmark suite. Radiometal--chelator
chemistry is likewise within your remit, as in Rudd \textit{et al.},
\textit{Chem. Sci.}, 2025, DOI: 10.1039/d4sc06969b. Every number and figure is
reproducible from the deposited dataset, geometries and code, which we release
as an open benchmark so the result can be overturned if it is wrong.

\paragraph{Suggested reviewers.}
We respectfully suggest referees spanning the three axes of the manuscript:
\begin{enumerate}
\item A coordination or radiometal chemist with experience of bifunctional
  chelator design and metal--chelate stability.
\item A computational chemist experienced in semi-empirical quantum chemistry
  (xTB family), implicit solvation and conformer sampling of metal complexes.
\item A researcher in chemistry-aware machine learning with expertise in
  validation design, data leakage and benchmark auditing.
\item A medicinal chemist or chemical biologist working on receptor-targeted
  conjugates and binding assays.
\item A statistician with experience of matched/paired designs and
  meta-analysis of heterogeneous assay data.
\end{enumerate}
Specific names will be supplied through the editorial system. Because the
dataset is drawn from published work by a small number of groups, we ask that
referees be selected from outside the laboratories whose data we re-analyse, and
we identify no other exclusions.

\paragraph{Conflicts and prior submission.}
The work has not been published or submitted elsewhere. There are no financial
or other conflicts to declare.

We thank you for your consideration and look forward to your editorial decision.

\closing{Sincerely yours,}

\end{letter}
\end{document}
